% ---------------------------------------------------------------------
%  CHƯƠNG 6 — KẾT LUẬN VÀ HƯỚNG PHÁT TRIỂN
% ---------------------------------------------------------------------
\chapter{Kết luận và hướng phát triển}
\label{ch:ketluan}

\begin{muctieu}
Chương này tổng kết đồ án, nêu bài học và hướng phát triển tiếp theo.
\end{muctieu}

\section{Kết luận}

Đồ án đã xây dựng thành công một ứng dụng mạng nhiều máy khách cho bài toán
bảng xếp hạng giải đấu: máy chủ TCP đa luộng nhận kết quả từ trọng tài, tính lại
bảng xếp hạng an toàn luồng và phát cập nhật tới mọi khán giả đang xem. Các kỹ
thuật bắt buộc của đề đều được dùng đúng chỗ: \code{Comparator.thenComparing}
cho xếp hạng, \code{record} cho trận, \code{Map} cho danh sách đội,
\code{CopyOnWriteArrayList} cho khách xem, \code{synchronized} cho trạng thái
dùng chung và socket đa luồng cho phục vụ nhiều máy khách. Hệ thống chạy đúng
trên dữ liệu thử, không mất trận khi tám trọng tài gửi đồng thời.

\section{Bài học rút ra}

Bài học lớn nhất là tách tầng dữ liệu khỏi tầng mạng. Việc đặt toàn bộ logic
xếp hạng trong \code{BangXepHang} và hàm \code{tinhHang} thuần dữ liệu giúp
kiểm thử nhanh và tin cậy trước khi viết một dòng socket nào, đúng tiến độ nhóm
đề A yêu cầu ở tuần~8. Bài học thứ hai là đồng bộ phải ở mức thao tác nghiệp vụ
chứ không phải mức trường: nếu không, hai luồng có thể cập nhật hai đội của
cùng một trận mà không thấy nhau. Nếu làm lại, em sẽ thêm cơ chế đóng khung
thông điệp ngay từ đầu thay vì sửa trong quá trình viết, vì nó ảnh hưởng tới
toàn bộ client.

\section{Phân công và lịch sử phiên bản}

Đồ án được quản lý bằng \code{git}, lịch sử commit trải đều từ tuần~5. Mỗi phần
(model, core, server, test, mở rộng) được commit riêng để dễ truy vết. Phiếu
phân công công việc và bản thiết kế hai trang tuần~5 được nộp cùng tiến độ. File
\code{git-log.txt} trong gói nộp là kết quả của \code{git log --stat}.

\section{Mở rộng đã làm}

Phần mở rộng đã hoàn thành là sinh lịch thi đấu vòng tròn một lượt tự động qua
lệnh \code{SCHEDULE <các đội>}, dùng phương pháp vòng tròn \en{(circle method)}:
cố định đội đầu, xoay các đội còn lại mỗi vòng. Nếu số đội lẻ thì thêm một đội
nghỉ \code{BYE} để lịch đối xứng. Kết quả trả về dưới dạng khung
\code{SCHEDULE <n>} kèm \(n\) dòng \code{<vòng> <độiNhà> <độiKhách>}. Phần này
đã được kiểm thử: bốn đội sinh đúng sáu trận, mỗi cặp đôi duy nhất, mỗi đội đá
ba trận; năm đội sinh đúng mười trận.

\section{Hướng phát triển}

Một số hướng phát triển tiếp theo:
\begin{itemize}[leftmargin=1.6em, itemsep=2pt]
  \item Thêm giao diện đồ hoạ bằng Swing cho client trọng tài và khán giả.
  \item Hỗ trợ UDP cho cập nhật bảng xếp hạng định kỳ, song song với TCP cho
        nhập kết quả.
  \item Thêm xác thực trọng tài để chỉ trọng tài mới được gửi \code{SUBMIT}.
  \item Hỗ trợ vòng tròn hai lượt \en{(double round-robin)} và xếp lịch theo sân
        nhà/sân khách cân bằng.
\end{itemize}
